\scanpage{69}
\mspage{18}

$W(6)$ (limites) m'apparaissent comme de nature plus spéciale, je ne songerais à imposer aucun d'eux pour l'\deleted{\unsure{introduction}} \inserted{usage} terminologique de ``localiseur fondamental''.

II) Perplexités sur les axiomes et la terminologie

\deleted{A propos de la notion de $W$-fibration, quand $W$ provient d'un dérivateur, on a présumé que \deleted{\unsure{les morphismes}} \inserted{une} $W$-fibration gauche, ou droite, est déjà une $W$-fibration}

Dans I et la plupart du temps, en travaillant avec \unsure{un dérivateur|des dérivateurs}, j'ai supposé $\boxed{W_{\mathbb D}\supset W_\infty}$, \deleted{\unsure{donc}} hypothèse qui m'a \unsure{parue|paru} essentielle (à tort ou à raison).
% slanted
\margin{à \uline{tort}, il s'avère \unsure{maintenant}.} On voit donc que si $W_{\mathbb D}$ satisfait à l'axiome des fibrations\inserted{$W(4)$}, même sous sa forme ``la plus faible $W(4\,\unsure{\mathrm d|\mathrm b}'')$, il est \uline{égal} à $W_\infty$. Cela montre \deleted{\unsure{que}} \inserted{bien} que cet axiome doit être vu comme \deleted{\unsure{prati}}quement une \uline{caractérisation} de $W_\infty$ (parmi les localiseurs fondamentaux qui contiennent $W_\infty$ -- et ces localiseurs me paraissent les plus importants, et peut-être les seuls vraiment utiles.)

La théorie des $\mathbb D$-fibrations de I, p. 111--\unsure{126|125}, a été faite dans l'esprit qu'on vient de dire, en supposant $W_{\mathbb D}\supset W_\infty$. \deleted{\unsure{Cette notion}} \deleted{\unsure{avait été}} \inserted{La notion de $\mathbb D$-fibration} \inserted{à laquelle \unsure{je m'étais finalement arrêté|j'étais finalement arrivé} était} par suite différente des notions envisagées \deleted{\unsure{de}} ici dans un esprit différent (\unsure{profa.|prop.} 4). Elle revient à mettre la condition a) (trop forte presque tjrs, \deleted{\unsure{qu}} sauf pour $W=W_\infty$), caractérisant les $W$-fibrations ultrafortes, par des $Y''\longrightarrow Y'$ où $Y''$, $Y'$ sont tous deux $W_\infty$-asphériques (ce qui est plus \uline{\unsure{exigent|exigeant}} que $W$-asphériques) ; \deleted{\unsure{c'est seulement}} qd $Y''$ et $Y'$ sont $W_\infty$-asphériques \inserted{(au lieu de $W$-asphériques)} qu'on demande que \deleted{la condition que} $X''\longrightarrow X'$ soit \deleted{\unsure{une}} dans $W$. On trouve dans

% archive sheet number: 69
